% ========================================================
% ANÁLISIS DEL CASO DE USO
% ========================================================
\section{Análisis del caso de uso}\label{sec:caso}

\subsection{Contexto y relación con el ODS 11}

El ODS~11 busca ciudades inclusivas, seguras, resilientes y sostenibles, y su meta~11.2 apunta a
sistemas de transporte seguros, asequibles, accesibles y sostenibles. Planificar ese transporte
requiere saber cuándo, cuánto y entre qué zonas se mueve la gente. Los registros de taxi responden
en parte a esa pregunta, porque cada viaje deja su hora de inicio y fin, su distancia y sus zonas
de origen y destino.

La \textit{Taxi and Limousine Commission} (TLC) de Nueva York publica estos registros como datos
abiertos \parencite{nyctlc}, con un diccionario de datos oficial y millones de viajes por mes. Ese
volumen hace que la concurrencia sea una necesidad del problema y no un requisito agregado.

\subsection{Pregunta del caso}

\begin{quote}
\textit{¿Qué tipos de viaje, según hora del día, día de la semana, duración y distancia, componen la
movilidad en taxi amarillo de Nueva York, y cómo se distribuyen entre las zonas de la ciudad?}
\end{quote}

Los tipos de viaje son los clusters de K-means. Su distribución por zona se obtiene después,
agregando las asignaciones por zona de origen y de destino, de modo que cada zona queda descrita
por su mezcla de tipos de viaje: un aeropuerto, por ejemplo, concentraría viajes largos a cualquier
hora.

\subsection{Por qué K-means}

K-means figura en la lista de modelos del enunciado y se ajusta a un problema sin variable objetivo.
Su costo, además, justifica la concurrencia. Cada iteración de Lloyd calcula la distancia de cada
viaje a cada centroide; con cerca de 2,8 millones de viajes y, por ejemplo, $K = 8$, son unos 22,6
millones de distancias en 6 dimensiones por iteración.

Dos propiedades hacen que se paralelice bien. La asignación de un viaje no depende de ningún otro,
así que los datos pueden repartirse en bloques entre goroutines, como en \textcite{nigro2022}. Y la
actualización de centroides se reduce a sumas y conteos por cluster, operaciones asociativas que
cada worker acumula por separado. Gracias a eso se puede comprobar que la versión concurrente
produce los mismos centroides que la secuencial.

\subsection{Alternativas consideradas}

El grupo evaluó primero un caso con datos de calidad del aire de SENAMHI, ligado al ODS~13, y optó
por NYC TLC por su volumen, su documentación y la trazabilidad de la fuente (PR~\#3 y PR~\#4 del
repositorio).

La primera versión de la propuesta agrupaba \emph{zonas}, con un punto por zona. Pero la ciudad
tiene 263 zonas identificables, y un K-means sobre 263 puntos no ejercita la concurrencia ni usa el
millón de registros que exige el enunciado. El modelo final agrupa \emph{viajes} y deja el análisis
por zona como agregación posterior. Esa decisión evita además un error de modelado: un
identificador de zona no es una coordenada, y usarlo en la distancia euclídea haría
\enquote{vecinas} a la zona 100 y a la 101 sin que lo sean.

\pendiente{confirmar que el caso y el modelo coinciden con la asignación del docente}

\subsection{Dataset seleccionado}

La Tabla~\ref{tab:dataset} resume el dataset. De sus 19 columnas, el caso usa siete
(Tabla~\ref{tab:columnas}); las demás, como proveedor, pasajeros, tipo de tarifa, medio de pago y
recargos, no intervienen en las reglas ni en las features.

\begin{table}[H]
\caption{Características del dataset}\label{tab:dataset}
\small
\begin{tabularx}{\textwidth}{@{}L{3.2cm}Y@{}}
\toprule
\textbf{Característica} & \textbf{Valor} \\
\midrule
Fuente & NYC Taxi \& Limousine Commission, \textit{TLC Trip Record Data} \\
Flota y periodo & Taxi amarillo (\textit{yellow}), enero de 2024 \\
Volumen & 2.964.624 viajes y 19 columnas en el archivo original \\
Formato & Parquet (49.961.641 bytes) y la tabla de zonas \texttt{taxi\_zone\_lookup.csv} \\
Licencia & Datos abiertos de la ciudad de Nueva York \\
\bottomrule
\end{tabularx}
\end{table}

\begin{table}[H]
\caption{Columnas del dataset que usa el caso}\label{tab:columnas}
\small
\begin{tabularx}{\textwidth}{@{}L{4.6cm}YL{4.2cm}@{}}
\toprule
\textbf{Columna} & \textbf{Contenido} & \textbf{Uso} \\
\midrule
\texttt{tpep\_pickup\_datetime}  & Inicio del viaje      & Reglas; hora, día y duración \\
\texttt{tpep\_dropoff\_datetime} & Fin del viaje         & Reglas; duración \\
\texttt{trip\_distance}          & Distancia en millas   & Reglas; feature \\
\texttt{PULocationID}            & Zona de origen        & Regla; agregación por zona \\
\texttt{DOLocationID}            & Zona de destino       & Regla; agregación por zona \\
\texttt{fare\_amount}            & Tarifa del taxímetro  & Regla de montos \\
\texttt{total\_amount}           & Monto total cobrado   & Regla de montos \\
\bottomrule
\end{tabularx}

{\footnotesize\textit{Nota.} Nombres del diccionario de datos de TLC. Los timestamps están en hora
local de Nueva York, sin zona horaria.}
\end{table}

\subsection{Diseño concurrente de referencia}\label{sec:diseno}

La implementación de PC2 partirá de un \emph{worker pool} de goroutines
(Figura~\ref{fig:workerpool}). Durante la asignación, los centroides son de solo lectura y cada
worker escribe únicamente en sus acumuladores privados, de modo que el bucle principal no comparte
estado mutable y no necesita \texttt{sync.Mutex}. La única sincronización ocurre al final de cada
iteración, en la barrera con \texttt{sync.WaitGroup}. El Algoritmo~\ref{alg:lloyd} detalla una
iteración.

\begin{figure}[H]
\caption{Worker pool para una iteración de Lloyd}\label{fig:workerpool}
\centering
\small
% Figura: worker pool de una iteración de Lloyd (sección 3.6).
\begin{tikzpicture}[
    font=\footnotesize,
    caja/.style={draw, rounded corners=1pt, align=center, inner sep=3pt, minimum height=9mm},
    flecha/.style={-Stealth, semithick}]
  \node[caja, text width=18mm] (datos) at (0.6, 0) {Viajes gold\\$N$ filas\\6 features};
  \node[caja, text width=44mm] (w1) at (5.2, 1.3) {Goroutine 1\\asigna y acumula $s_{1j}, n_{1j}$};
  \node (dots) at (5.2, 0.1) {$\vdots$};
  \node[caja, text width=44mm] (wP) at (5.2, -1.3) {Goroutine $P$\\asigna y acumula $s_{Pj}, n_{Pj}$};
  \node[draw, fill=black!15, minimum width=2.5mm, minimum height=36mm, inner sep=0] (barrera) at (9.0, 0) {};
  \node[above=1mm of barrera] {\texttt{sync.WaitGroup}};
  \node[caja, text width=40mm] (reduce) at (12.3, 0) {Reducción\\$\mu_j = \sum_w s_{wj} \,/\, \sum_w n_{wj}$};
  \draw[flecha] (datos.east) -- node[pos=0.55, above, sloped, font=\scriptsize] {$B_1$} (w1.west);
  \draw[flecha] (datos.east) -- node[pos=0.55, below, sloped, font=\scriptsize] {$B_P$} (wP.west);
  \draw[flecha] (w1.east) -- (w1.east -| barrera.west);
  \draw[flecha] (wP.east) -- (wP.east -| barrera.west);
  \draw[flecha] (barrera.east) -- (reduce.west);
  \draw[flecha, dashed] (reduce.south) -- ++(0, -2.2)
    -| node[pos=0.25, below] {nuevos centroides, de solo lectura en la siguiente iteración} (wP.south);
\end{tikzpicture}

\end{figure}

\begin{algorithm}[tbp]
\caption{Una iteración de Lloyd con worker pool y acumuladores privados}\label{alg:lloyd}
\begin{algorithmic}[1]
\Require viajes $X = \{x_1, \dots, x_N\} \subset \mathbb{R}^6$, centroides $\mu_1, \dots, \mu_K$,
         $P$ workers
\Ensure centroides actualizados e inercia $J$
\State Partir $X$ en $P$ bloques contiguos $B_1, \dots, B_P$ de tamaño similar
\ForAll{worker $w \in \{1, \dots, P\}$} \Comment{una goroutine; $\mu$ es de solo lectura}
  \State $s_{wj} \gets \mathbf{0}$, \ $n_{wj} \gets 0$ para $j = 1, \dots, K$; \ $J_w \gets 0$
  \ForAll{$x \in B_w$}
    \State $j^{*} \gets \arg\min_j \lVert x - \mu_j \rVert^2$
    \State $s_{wj^{*}} \gets s_{wj^{*}} + x$; \ $n_{wj^{*}} \gets n_{wj^{*}} + 1$;
           \ $J_w \gets J_w + \lVert x - \mu_{j^{*}} \rVert^2$
  \EndFor
\EndFor
\State Esperar a los $P$ workers \Comment{barrera: \texttt{sync.WaitGroup}}
\ForAll{cluster $j$}
  \If{$\sum_w n_{wj} > 0$}
    \State $\mu_j \gets \sum_w s_{wj} \,/\, \sum_w n_{wj}$ \Comment{sumas y conteos, nunca promedios de promedios}
  \EndIf
\EndFor
\State \Return $\mu$, \ $J = \sum_w J_w$
\end{algorithmic}
\end{algorithm}

La revisión bibliográfica fija cuatro decisiones de la evaluación. Primero, la versión concurrente
suma en otro orden que la secuencial, y en punto flotante eso altera los últimos dígitos; por eso
compararemos centroides e inercia con una tolerancia explícita \parencite{he2022}, partiendo de la
misma inicialización y el mismo criterio de parada \parencite{nigro2022}. Segundo, el tamaño de
bloque será una variable del experimento junto con el número de workers, porque el reparto afecta
el rendimiento tanto como la cantidad de hilos \parencite{kwedlo2021,kwedlo2024}. Tercero, el
speedup ($T_{\text{secuencial}} / T_{\text{concurrente}}$) se medirá contra la versión secuencial
del mismo algoritmo, con media recortada de varias ejecuciones, porque la cifra depende del
denominador \parencite{bellavita2025}.

Cuarto, la parte serial pone techo al speedup. Según la ley de Amdahl \parencite{amdahl1967}, si
una fracción $f$ del trabajo es paralelizable, con $P$ goroutines el speedup es
\[
  S(P) = \frac{1}{(1-f) + f/P}, \qquad S(P) < \frac{1}{1-f}.
\]
En nuestro diseño son seriales la lectura del CSV de gold, la preparación de los bloques, la
barrera y la actualización de centroides. Si, por ejemplo, el 14\,\% del tiempo total fuera
serial, con 8 goroutines el speedup no pasaría de unas 4 veces, y ni con infinitas llegaría a 7.
Por eso mediremos por separado la lectura y el tiempo por iteración, para estimar $f$. Además del
escalamiento fuerte (el mismo problema con más goroutines), reportaremos el escalamiento débil
(más viajes en el mismo tiempo al sumar goroutines), que describe la ley de Gustafson
\parencite{gustafson1988}. La memoria, en cambio, no es el límite: las 6 features de 2,8 millones
de viajes ocupan unos 136~MB en doble precisión.

Como contraste, PC2 medirá una variante con acumuladores compartidos protegidos por
\texttt{sync.Mutex}, para cuantificar el costo de la contención. La barrera se modelará en Promela
para verificar que no hay condiciones de carrera.
